\begin{lemma}\label{lem:nonic-cont-odd-cuspidal}
Fix $d\ge1$. Let $X/\mathbb Q$ be a modular curve with a rational cusp $c_0$,
rank-zero rational Jacobian and rational gonality greater than $d$.
Suppose $X$ has good reduction at an odd prime $p$, with finite etale
cusp scheme, and its cusps are defined over a number field unramified at $p$.
No noncuspidal effective rational divisor of degree $d$ on $X$ can
reduce entirely to cusps.
\end{lemma}
\begin{proof}
Suppose $D$ has entirely cuspidal reduction. Over a number field $L$
unramified at $p$ and containing all cusp coordinates, lift its reduction
to an effective cuspidal divisor $C$ of degree $d$. Both
$[D-dc_0]$ and $[C-dc_0]$ are torsion, by rank zero and Manin--Drinfeld,
respectively. Torsion reduction at an unramified odd prime is injective,
so $D\sim C$. Since $D\ne C$, base change gives
$h^0(X,D)=h^0(X_L,D)\ge2$, contrary to the gonality assumption.
\end{proof}

\begin{lemma}\label{lem:nonic-cont-hecke}
Fix $d\ge1$. Let $X/\mathbb Q$ have rank-zero Jacobian and a rational cusp $c_0$.
Let $p$ be a prime of good reduction, and write the reduction of an
effective rational divisor $D$ of degree $d$ as $D'+C'$, with $C'$
cuspidal. For an odd prime $q$ coprime to the level and $p$, put
$A_q=T_q-1-q\langle q\rangle$. If cyclotomic $q$ fixes every
Frobenius-$p$ cusp orbit which can occur in $C'$, then $A_qD'$ is principal.
\end{lemma}
\begin{proof}
Injectivity of rational torsion reduction at $q$ and Eichler--Shimura
give $A_q\operatorname{Jac}(X)(\mathbb Q)=0$.
Moreover, $A_qc_0=0$ as a divisor. Hence $A_qD$ is principal.
On cusps we have
$T_qC'=\sigma_q C'+q\langle q\rangle\sigma_q^{-1}C'$.
The hypothesis gives $A_qC'=0$ as a divisor, proving the assertion.
\end{proof}

For $X_1(N)$ the program determines the admissible $q$ using
the cusp description of \cite[Lemma~2.7]{DEHMZ}. It checks that
the predicted finite-field cusp degrees agree with the model.
Every integer partition of each noncuspidal degree at most $d$
is included, and repetitions of places are allowed. Diamond operators
commute with $A_q$ and preserve the admissible cusp complements, so it
suffices to take one diamond representative for the first support place
and all choices for the other support places.

\begin{lemma}\label{lem:nonic-cont-quotient-cusps}
Let $\pi:X\longrightarrow Y$ be a finite morphism of smooth projective
curves over $\mathbb Q$, taking cusps to cusps, and let $c_0$ be a
rational cusp on $Y$. Suppose $\operatorname{Jac}(Y)(\mathbb Q)$ is
generated by rational cuspidal divisor classes. Let $p$ be a common
prime of good reduction and let $H_p\subseteq\operatorname{Jac}(Y)(\mathbb F_p)$
be the image of these global cuspidal classes. Write $q_p$ for the
quotient by $H_p$, and put
\[
 S_k=\{q_p([B-kc_0]): B\text{ effective cuspidal on }Y_{\mathbb F_p},
 \ \deg B=k\}.
\]
If an effective rational divisor of degree $d$ on $X$ reduces to
$D'+B'$, with $B'$ cuspidal and $\deg D'=d'$, then
\[
 -q_p([\pi_*D'-d'c_0])\in S_{d-d'}.
\]
\end{lemma}
\begin{proof}
The global class $[\pi_*D-dc_0]$ is rational, hence cuspidal by the
hypothesis. Reduction and proper pushforward commute. Both operations
preserve the degree of a divisor; in particular a closed point $P$
with image $Q$ pushes forward to $[\kappa(P):\kappa(Q)]Q$.
The image of $[\pi_*D-dc_0]$ in the quotient is therefore zero, and
$\pi_*B'$ is an effective cuspidal divisor of degree $d-d'$.
This gives the asserted membership. Allowing every finite-field
cuspidal complement on $Y$ gives a necessary condition, even when
some of them do not lift to global divisors.
\end{proof}
